\documentclass{stp}

\newcommand{\en}[1]{\textit{#1}}

\begin{document}

\stptitle{
  вид        = контрольная,
  номер      = 1,
  тема       = {Серверная разработка},
  дисциплина = {ИТиВП},
  студент    = {А. С. Образцов},
  группа     = {320601},
  преподаватель = {Н. В. Хаджинова},
}

\section{Цель работы}

Целью контрольной работы является разработка \en{REST API} на платформе
\en{Node.js} с аутентификацией по \en{JWT}, хранением данных в
\en{PostgreSQL} и полным набором \en{CRUD}"=операций для сущности,
соответствующей теме курсового проекта. Тема курсового проекта --- сервис
планирования задач по методу \en{GTD} (\en{Getting Things Done}) с
контекстами и напоминаниями (вариант~27).

Метод \en{GTD} предполагает, что все дела человека сначала собираются во
входящие, а затем разбираются по спискам. Поэтому основной сущностью
сервиса выбрана задача (\en{Task}): она хранит заголовок, заметки, статус
(\en{inbox}, \en{next}, \en{waiting}, \en{scheduled}, \en{someday},
\en{done}), срок выполнения \en{dueDate}, время напоминания \en{remindAt},
контекст \en{contextId} и владельца \en{ownerId}. Помимо стандартных полей
\en{id}, \en{createdAt} и \en{updatedAt} сущность содержит семь собственных
полей, что превышает минимум в четыре поля. Поле \en{status} ограничено
перечислением, поле \en{title} и обе даты проходят валидацию, а уникальность
обеспечивается полем \en{email} модели пользователя.

\section{Ход работы}\label{sec:work}

\subsection{Выбор технологий и структура проекта}

Сервис написан на \en{Node.js} с фреймворком \en{Express~5}. Данные хранятся
в \en{PostgreSQL~16}, запущенной в контейнере \en{Docker}, доступ к базе
выполняется через \en{ORM Sequelize}. Пароли хешируются библиотекой
\en{bcrypt}, токены подписываются библиотекой \en{jsonwebtoken}, входные
данные проверяются схемами \en{Joi}. Для защиты используются \en{helmet},
\en{cors} и \en{express-rate-limit}. Выбор \en{PostgreSQL} вместо
\en{MongoDB} объясняется структурой данных: у задач, контекстов и
пользователей есть чёткие связи и ограничения целостности, которые
реляционная база проверяет сама.

Код разделён на слои: маршруты (\en{routes}) описывают цепочку
\en{middleware}, контроллеры (\en{controllers}) читают проверенные данные
запроса и вызывают сервисы (\en{services}), в которых сосредоточены
бизнес-логика и проверка прав доступа, а модели (\en{models}) описывают
таблицы и связи. Настройки хранятся в файле \en{.env}, образец
\en{.env.example} лежит в репозитории.

Исходный код размещён в репозитории
\en{https://github.com/\allowbreak{}Alexandr-Obraztsov/\allowbreak{}gtd-task-service},
ветка \en{kr21}. Коллекция \en{Postman} находится в файле
\en{postman/\allowbreak{}gtd-secure-api.postman\_collection.json}, документация
--- в файлах \en{API.md} и \en{docs/\allowbreak{}openapi.json}, инструкция по
запуску --- в \en{README.md}.

\subsection{Схема базы данных}

База данных состоит из четырёх таблиц. Таблица \en{users} хранит учётные
записи: уникальный \en{email}, хеш пароля, роль (\en{user}, \en{moderator},
\en{admin}), счётчик неудачных входов и время окончания блокировки. Таблица
\en{refresh\_tokens} хранит хеши выданных \en{refresh}-токенов. Таблицы
\en{contexts} и \en{tasks} хранят контексты и задачи, у каждой записи есть
владелец \en{ownerId}. Схема базы данных представлена на
рисунке~\ref{fig:er}.

\stpimage[160mm]{img/er.png}{Схема базы данных сервиса}{fig:er}

При удалении пользователя каскадно удаляются его задачи, контексты и
токены. При удалении контекста задачи не удаляются: поле \en{contextId} у
них становится пустым.

\subsection{Модели}

Модель пользователя содержит уникальный \en{email} с проверкой формата и
роль из перечисления. Метод \en{toJSON} возвращает только открытые поля,
поэтому хеш пароля не может попасть в ответ даже по ошибке в контроллере:

\begin{stplisting}
const PUBLIC_FIELDS = Object.freeze(
  ['id', 'email', 'role', 'createdAt', 'updatedAt']);

class User extends Model {
  toJSON() {
    const values = this.get({ plain: true });
    return Object.fromEntries(
      PUBLIC_FIELDS.map((field) => [field, values[field]]));
  }
}

User.init({
  id: { type: DataTypes.INTEGER, autoIncrement: true,
    primaryKey: true },
  email: { type: DataTypes.STRING(254), allowNull: false,
    unique: true, validate: { isEmail: true } },
  passwordHash: { type: DataTypes.STRING(100),
    allowNull: false },
  role: { type: DataTypes.ENUM(...ROLE_VALUES),
    allowNull: false, defaultValue: ROLES.USER },
}, { sequelize, modelName: 'User', tableName: 'users' });
\end{stplisting}

Модель задачи ограничивает статус перечислением \en{GTD} и задаёт индексы
для выборки по владельцу и статусу и для поиска наступивших напоминаний:

\begin{stplisting}
Task.init({
  id: { type: DataTypes.INTEGER, autoIncrement: true,
    primaryKey: true },
  title: { type: DataTypes.STRING(200), allowNull: false },
  notes: { type: DataTypes.TEXT, allowNull: false,
    defaultValue: '' },
  status: { type: DataTypes.ENUM(...TASK_STATUS_VALUES),
    allowNull: false, defaultValue: TASK_STATUS.INBOX },
  dueDate: { type: DataTypes.DATE, allowNull: true },
  remindAt: { type: DataTypes.DATE, allowNull: true },
  contextId: { type: DataTypes.INTEGER, allowNull: true },
  ownerId: { type: DataTypes.INTEGER, allowNull: false },
}, {
  sequelize, modelName: 'Task', tableName: 'tasks',
  indexes: [{ fields: ['ownerId', 'status'] },
    { fields: ['remindAt'] }],
});
\end{stplisting}

\subsection{Маршруты, роли и права доступа}

Маршруты задач требуют аутентификации для всех операций, а каждый маршрут
имеет собственную схему проверки входных данных. Маршрут
\en{GET /tasks} принимает необязательные фильтры \en{status},
\en{contextId} и \en{dueBefore}, а также постраничный вывод --- это
дополнительный эндпоинт с фильтрацией, которого требует задание. Для
\en{GET /tasks/reminders/due} возвращаются задачи пользователя, время
напоминания которых уже наступило:

\begin{stplisting}
const router = Router();

router.use(authenticate);

router.get('/', validate(schemas.list), controller.list);
router.post('/', validate(schemas.create), controller.create);
router.get('/reminders/due', validate(schemas.dueReminders),
  controller.dueReminders);
router.get('/:id', validate(schemas.getOne), controller.getOne);
router.put('/:id', validate(schemas.replace),
  controller.replace);
router.patch('/:id', validate(schemas.update),
  controller.update);
router.delete('/:id', validate(schemas.remove),
  controller.remove);
\end{stplisting}

Права доступа заданы заданием контрольной работы: создавать задачи может
любая аутентифицированная роль (\en{user} и выше), обновлять --- автор
задачи или администратор, удалять --- только администратор. Проверка
выполняется в сервисе, поскольку для неё нужна сама запись из базы данных.
Администратору разрешено обновление, остальные роли обновляют только свои
задачи. Удаление проверяет роль до обращения к базе, поэтому ответ не
раскрывает существование чужой задачи:

\begin{stplisting}
const TASK_ACCESS_POLICY = Object.freeze({
  ...PRIVILEGED_ROLE_BY_ACTION,
  [ACTIONS.UPDATE]: ROLES.ADMIN,
});

function replaceTask(actor, id, input) {
  return updateTask(actor, id,
    { ...FULL_UPDATE_DEFAULTS, ...input });
}

async function deleteTask(actor, id) {
  if (!hasRole(actor.role, ROLES.ADMIN)) {
    throw new ForbiddenError(
      'Удалять задачи может только администратор',
      'ADMIN_ONLY');
  }
  const task = await loadTask(actor, id, ACTIONS.DELETE);
  await task.destroy();
}
\end{stplisting}

Метод \en{PUT} выполняет полную замену: обязателен только заголовок, а не
переданные поля (\en{notes}, \en{status}, \en{dueDate}, \en{remindAt},
\en{contextId}) сбрасываются к значениям по умолчанию. Метод \en{PATCH}
меняет только переданные поля.

\subsection{Валидация входных данных}

Для каждого источника данных каждого маршрута (\en{params}, \en{query},
\en{body}) описана схема \en{Joi}, неизвестные поля запрещены. Схема
создания задачи требует заголовок, ограничивает его длину, проверяет
статус по перечислению и приводит даты к формату \en{ISO~8601} в
допустимом диапазоне; схема замены повторяет её:

\begin{stplisting}
const title = Joi.string().trim().min(1).max(200);
const notes = Joi.string().trim().max(2000).allow('');
const status = Joi.string().valid(...TASK_STATUS_VALUES);
const boundedDate = Joi.date().iso()
  .min('2000-01-01T00:00:00Z').max('2100-12-31T23:59:59Z');

const create = {
  body: Joi.object({
    title: title.required(),
    notes,
    status,
    dueDate: boundedDate.allow(null),
    remindAt: boundedDate.allow(null),
    contextId: id.allow(null),
  }),
};

const replace = { params: idParams, body: create.body };
\end{stplisting}

Для регистрации пароль должен содержать не менее восьми символов, букву,
цифру и спецсимвол, а \en{email} приводится к нижнему регистру и
проверяется на формат. Кроме схем, действует ограничение тела запроса
размером 10~КБ и приём только формата \en{application/json}.

\subsection{Аутентификация}

При регистрации пароль хешируется алгоритмом \en{bcrypt} с 12 раундами, в
базу попадает только хеш. Уникальность \en{email} проверяется запросом к
базе и дополнительно гарантируется ограничением \en{UNIQUE} на столбец;
при повторной регистрации сервер отвечает статусом 409:

\begin{stplisting}
async function registerUser({ email, password }) {
  const existing = await User.findOne({ where: { email } });
  if (existing) {
    throw new ConflictError(
      'Пользователь с таким email уже существует');
  }
  const passwordHash = await hashPassword(password);
  const user = await User.create(
    { email, passwordHash, role: ROLES.USER },
    { fields: ['email', 'passwordHash', 'role'] });
  return user;
}
\end{stplisting}

При входе пароль сравнивается с хешем, после чего выдаётся пара токенов:
короткоживущий \en{access}"=токен (15 минут) возвращается в теле ответа, а
\en{refresh}"=токен --- в \en{httpOnly cookie}. \en{Middleware authenticate}
проверяет формат заголовка \en{Authorization}, подпись и срок действия
токена, а затем загружает пользователя из базы. Роль берётся из базы, а не
из токена, поэтому понижение роли действует сразу:

\begin{stplisting}
async function authenticate(req, res, next) {
  const payload = decodeAccessToken(req,
    extractBearerToken(req));
  const user = await User.findByPk(Number(payload.sub));
  if (!user) {
    throw new UnauthorizedError('Пользователь не найден',
      'TOKEN_INVALID');
  }
  req.actor = { id: user.id, email: user.email,
    role: user.role, ip: req.ip };
  next();
}
\end{stplisting}

\subsection{Проверка регистрации и входа}

Все запросы выполнялись к работающему серверу, ответы на рисунках
получены от него без изменений; значения токенов и паролей сокращены или
скрыты. Результат регистрации нового пользователя представлен на
рисунке~\ref{fig:register}. Сервер вернул статус 201 и открытые поля
пользователя, роль назначена автоматически.

\stpimage[150mm]{img/register.png}{Регистрация нового пользователя}{fig:register}

Повторная регистрация с тем же \en{email} отклонена: ответ приведён на
рисунке~\ref{fig:registerdup}. Это подтверждает уникальность адреса.

\stpimage[150mm]{img/registerDup.png}{Повторная регистрация с занятым \en{email}}{fig:registerdup}

Запрос с некорректными данными показан на рисунке~\ref{fig:registerbad}.
Сервер вернул статус 400~с кодом \en{VALIDATION\_ERROR} и списком ошибок по
каждому полю.

\stpimage[125mm]{img/registerBad.png}{Ошибки валидации при регистрации}{fig:registerbad}

Вход с верными данными представлен на рисунке~\ref{fig:login}. В ответе
приходят данные пользователя и \en{JWT} для заголовка \en{Authorization}.

\stpimage[150mm]{img/login.png}{Вход и получение \en{JWT}}{fig:login}

Вход с неверным паролем завершается статусом 401 и общим сообщением, не
выдающим существование учётной записи (рисунок~\ref{fig:loginbad}).

\stpimage[150mm]{img/loginBad.png}{Вход с неверным паролем}{fig:loginbad}

Доступ к защищённому маршруту без токена и с поддельным токеном показан на
рисунках~\ref{fig:notoken} и~\ref{fig:badtoken}: в обоих случаях сервер
возвращает статус 401~с разными кодами ошибки.

\stpimage[150mm]{img/noToken.png}{Запрос к защищённому маршруту без токена}{fig:notoken}

\stpimage[150mm]{img/badToken.png}{Запрос с поддельным токеном}{fig:badtoken}

\subsection{Проверка \en{CRUD}-операций}

Создание задачи с токеном пользователя показано на
рисунке~\ref{fig:create}. Сервер вернул статус 201, присвоил задаче
идентификатор и владельца, а даты привёл к единому формату.

\stpimage[140mm]{img/createTask.png}{Создание задачи}{fig:create}

Получение списка с фильтрацией по статусу и контексту (дополнительный
эндпоинт) приведено на рисунке~\ref{fig:list}. В ответ попадают только
задачи, подходящие под оба условия, а вместе с ними --- общее число
записей и параметры страницы.

\stpimage[135mm]{img/listTasks.png}{Список задач с фильтрацией}{fig:list}

Получение одной задачи по идентификатору представлено на
рисунке~\ref{fig:get}.

\stpimage[140mm]{img/getTask.png}{Получение задачи по идентификатору}{fig:get}

Полная замена задачи методом \en{PUT} показана на рисунке~\ref{fig:put}.
Не переданные поля сброшены к значениям по умолчанию.

\stpimage[140mm]{img/putTask.png}{Полная замена задачи методом \en{PUT}}{fig:put}

Ошибки валидации при создании и при замене задачи приведены на
рисунках~\ref{fig:createbad} и~\ref{fig:putbad}: пустой заголовок и
недопустимый статус, а также отсутствие обязательного поля \en{title}
отклоняются статусом 400.

\stpimage[135mm]{img/createBad.png}{Создание задачи с недопустимыми данными}{fig:createbad}

\stpimage[135mm]{img/putBad.png}{Замена задачи без заголовка}{fig:putbad}

Права на удаление проверены двумя запросами. Попытка удалить задачу от
имени пользователя с ролью \en{user} показана на
рисунке~\ref{fig:deleteuser}: сервер вернул статус 403~с кодом
\en{ADMIN\_ONLY}, даже несмотря на то что пользователь является владельцем.

\stpimage[150mm]{img/deleteUser.png}{Попытка удаления от имени \en{user}}{fig:deleteuser}

Удаление той же задачи администратором представлено на
рисунке~\ref{fig:deleteadmin}: сервер вернул статус 204 без тела ответа.

\stpimage[150mm]{img/deleteAdmin.png}{Удаление задачи администратором}{fig:deleteadmin}

Повторное обращение к удалённой задаче завершилось статусом 404
(рисунок~\ref{fig:getdeleted}).

\stpimage[150mm]{img/getDeleted.png}{Запрос удалённой задачи}{fig:getdeleted}

\subsection{Документация и автоматическое тестирование}

\en{API} описан в файле \en{API.md} и в спецификации \en{OpenAPI~3.0},
которая отдаётся интерактивной страницей \en{Swagger UI} по адресу
\en{/api-docs}. Раздел задач из \en{Swagger UI} приведён на
рисунке~\ref{fig:swagger}: видны все маршруты, включая \en{PUT} и
\en{DELETE} с пометкой <<только \en{admin}>>.

\stpimage[150mm]{img/swagger.png}{Раздел задач в \en{Swagger UI}}{fig:swagger}

Для повторяемой проверки подготовлена коллекция \en{Postman} из 129
запросов и 196 проверок: регистрация, вход, \en{CRUD}, ошибки валидации,
доступ без токена, чужие ресурсы и роли. Прогон коллекции через
\en{Newman} завершился без ошибок, результат показан на
рисунке~\ref{fig:newman}.

\stpimage[150mm]{img/newman.png}{Прогон коллекции \en{Postman} через \en{Newman}}{fig:newman}

\section{Ответы на контрольные вопросы}

\begin{stpnum}
  \item Выбрана сущность <<задача>> (\en{Task}): сервис планирования по
    методу \en{GTD} --- тема курсового проекта, а задача является его
    центральным объектом. Сущность содержит заголовок, заметки, статус,
    срок, время напоминания, контекст и владельца.
  \item Использованы \en{Node.js} и \en{Express~5} как привычный и
    быстрый стек для \en{REST API}, \en{PostgreSQL} и \en{Sequelize} для
    хранения связанных данных с проверкой целостности, \en{jsonwebtoken} и
    \en{bcrypt} для аутентификации, \en{Joi} для валидации, \en{helmet},
    \en{cors} и \en{express-rate-limit} для защиты.
  \item Пользователь регистрируется запросом \en{POST /auth/register}:
    пароль хешируется, запись сохраняется с ролью \en{user}. При входе
    \en{POST /auth/login} пароль сверяется с хешем, после чего выдаются
    \en{access}"=токен в теле ответа и \en{refresh}"=токен в \en{httpOnly
    cookie}. Защищённые маршруты проверяют заголовок
    \en{Authorization: Bearer}.
  \item Пароли хешируются алгоритмом \en{bcrypt} с 12 раундами; соль
    генерируется автоматически и хранится в хеше. Это важно, потому что
    при утечке базы злоумышленник получает не пароли, а хеши, подбор
    которых требует много времени и ресурсов.
  \item Реализованы роли \en{user}, \en{moderator} и \en{admin}.
    Пользователь создаёт задачи и работает со своими. Модератор
    дополнительно читает чужие задачи. Администратор дополнительно
    изменяет любые задачи, удаляет задачи и управляет пользователями.
    Удалять задачи может только администратор.
  \item Перед созданием пользователя выполняется поиск по \en{email}, при
    совпадении возвращается статус 409. Дополнительно на столбец
    наложено ограничение \en{UNIQUE}, которое защищает от гонки двух
    одновременных регистраций. Адрес предварительно приводится к нижнему
    регистру.
  \item \en{GET} --- получение списка и отдельной задачи, \en{POST} ---
    регистрация, вход и создание задач, \en{PUT} --- полная замена задачи,
    \en{PATCH} --- частичное изменение, \en{DELETE} --- удаление задачи
    администратором.
  \item Каждый маршрут имеет схему \en{Joi}. При ошибке возвращается
    статус 400 и объект с кодом \en{VALIDATION\_ERROR} и списком
    нарушений. Пример ответа:
    \en{\{ "error": \{ "code": "VALIDATION\_ERROR", "message":
    "Некорректные данные запроса", "details": [ \{ "location": "body",
    "field": "email", "message": "должно быть корректным email" \} ] \}
    \}}.
  \item В коде нет сырого \en{SQL}: все запросы строятся методами
    \en{Sequelize}, которые передают значения как параметры, а не
    подставляют их в текст запроса. Идентификаторы дополнительно
    приводятся \en{Joi} к целым числам.
  \item Применены \en{helmet} (строгая \en{CSP}, \en{nosniff}, запрет
    встраивания), ограничение частоты запросов (глобально 300 за 15 минут,
    для маршрутов \en{/auth/*} --- 20 неуспешных), \en{CORS} с белым
    списком источников, лимит тела запроса 10~КБ, приём только
    \en{application/json}, санитизация \en{HTML} в строках и блокировка
    учётной записи после пяти неудачных входов.
  \item Использованы \en{Swagger UI} по спецификации \en{OpenAPI~3.0} и
    файл \en{API.md}. Спецификация машиночитаема, по ней можно сразу
    пробовать запросы, а \en{API.md} удобно читать в репозитории.
  \item Использованы коллекция \en{Postman} из 129 запросов с 196
    проверками, запуск через \en{Newman} из консоли и ручные запросы к
    работающему серверу, результаты которых приведены на рисунках.
  \item Создание --- 201 \en{Created}, успешное обновление --- 200
    \en{OK}, удаление --- 204 \en{No Content}. Ошибочные случаи: 400, 401,
    403, 404 и 409.
  \item Секреты (\en{JWT\_SECRET}, \en{JWT\_REFRESH\_SECRET}, строка
    подключения к базе) хранятся в файле \en{.env}, который не попадает в
    репозиторий. В репозитории лежит только \en{.env.example} с
    заглушками. При запуске сервер не принимает секреты короче 32
    символов.
  \item \en{CORS} (\en{Cross-Origin Resource Sharing}) --- механизм, по
    которому браузер разрешает странице с одного источника обращаться к
    \en{API} другого источника. Сервер принимает запросы только с
    источников из переменной \en{CORS\_ORIGINS} (адреса разработки
    клиента \en{http://localhost:5173} и \en{http://localhost:3000}),
    запросы с чужих источников отклоняются.
  \item Планируется подключить клиент на \en{React}, добавить рассылку
    напоминаний по расписанию, повторяющиеся задачи, общий доступ к
    спискам и развёртывание на облачной платформе.
\end{stpnum}

\stpunnumbered{Заключение}

В ходе контрольной работы разработан \en{REST API} сервиса планирования
задач по методу \en{GTD} на \en{Node.js} и \en{Express}. Реализованы
регистрация и вход с \en{JWT}, хранение данных в \en{PostgreSQL} через
\en{Sequelize}, полный набор \en{CRUD}"=операций для сущности <<задача>>
с правами по ролям: создание --- любая аутентифицированная роль, обновление
--- автор или администратор, удаление --- только администратор.

Все входные данные проверяются схемами \en{Joi}, ошибки возвращаются в
едином формате с кодами, результаты проверок приведены в
разделе~\ref{sec:work}. Добавлены фильтрация списка задач и выборка
сработавших напоминаний, подготовлены документация в \en{Swagger UI},
файле \en{API.md} и коллекция \en{Postman}, прогон которой через
\en{Newman} прошёл без ошибок.

Код размещён в репозитории
\en{https://github.com/\allowbreak{}Alexandr-Obraztsov/\allowbreak{}gtd-task-service}
в ветке \en{kr21}, коллекция \en{Postman} --- по адресу
\en{https://github.com/\allowbreak{}Alexandr-Obraztsov/\allowbreak{}gtd-task-service/\allowbreak{}blob/\allowbreak{}kr21/\allowbreak{}postman/\allowbreak{}gtd-secure-api.postman\_collection.json}.
Разработанный \en{API} служит серверной частью курсового проекта и будет
дополнен клиентским приложением.

\end{document}
